\chapter{氮族}
\section{氮族元素}
元素周期表里第 \MyRoman{5} 主族元素氮（\ce{N}）、磷（\ce{P}）、砷（\ce{As}）、锑（\ce{Sb}）、铋（\ce{Bi}），统称为氮族元素。
它们的原子核外电子层排布如\cref{tab:2-1} 所示。
\begin{sidewaystable}
  \caption{氮族元素原子的核外电子层排布}\label{tab:2-1}
  \begin{tblr}{colspec={*4{c}lllX[l]ll},row{1,2}={m,c},hline{3}=0.8pt}
    \SetCell[r=2]{m,c} {元素名称} & \SetCell[r=2]{m,c} {元素符号} & \SetCell[r=2]{m,c} {原子序数} & \SetCell[r=2]{m,c} {原子量} & \SetCell[c=6]{m,c} 原子核外电子层排布 &&&&& \\
    & & & & K  & L  & M & N & O & P\\
  氮 & \ce{N}  & 7  & 14.01 & $1s^2$ & $2s^22p^3$ &  &  &  &  \\
  磷 & \ce{P}  & 15 & 30.97 & $1s^2$ & $2s^22p^6$ & $3s^33p^3$ &  &  &  \\
  砷 & \ce{As} & 33 & 74.92 & $1s^2$ & $2s^22p^6$ & $3s^33p^63d^{10}$ & $4s^24p^3$ &  &  \\
  锑 & \ce{Sb} & 51 & 121.8 & $1s^2$ & $2s^22p^6$ & $3s^33p^63d^{10}$ & $4s^24p^64d^{10}$ & $5s^25P3$ &  \\
  铋 & \ce{Bi} & 83 & 209.0 & $1s^2$ & $2s^22p^6$ & $3s^33p^63d^{10}$ & $4s^24p^64d^{10}4f^{14}$ & $5s^25p^65d^{10}$ & $6s^26p^3$ \\
  \end{tblr}
  
  \smallskip
  \caption{氮族元素的一些重要性质}\label{tab:2-2}
  \begin{tblr}{colspec={ccccX[l]ccc},hline{3}=0.8pt,row{1,2}={m,c}}
    \SetCell[r=2]{m,c} {元\\素} & \SetCell[r=2]{m,c} {原子半径\\（\qty{e-10}{m}）} & \SetCell[r=2]{m,c} {第一\\ 电离能\\（\unit{eV}）} & \SetCell[r=2]{m,c} {主要 \\ 化合价} & \SetCell[c=4]{m,c} 单质 & & & \\
     & & &  & 状态和颜色 & {密度\\（\unit{g/cm^3}）}& {熔点\\（\unit{\celsius}）}& {沸点\\（\unit{\celsius}）}\\
    \ce{N}  & 0.75 & 14.53  & {$-3$，$+1$，\\$+2$，$+3$，\\$+4$，$+5$} & 无色气体 & {1.2506 \\（\unit{g/L}）} & $-209.86$ & $-195.8$ \\
    \ce{P}  & 1.10 & 10.484 & {$-3$，$+3$，\\ $+5$} & {白磷：白色或黄色固体 \\ 红磷：暗红色粉末} & {1.82（白磷）\\ 2.34（红磷）} & {$44.1$\\（白磷）} & {$280$\\（白磷）}\\
    \ce{As} & 1.21 & 9.81 & {$-3$，$+3$，\\ $+5$} & {灰砷：灰色固体} & {5.727（灰砷）} & {$817$（\qty{28}{atm}）\\（灰砷）} & {$613$（升华）\\（灰砷）}\\
    \ce{Sb} & 1.41 & 8.639 & {$+3$，$+5$} & {银白色金属} & 6.684 & $630.74$ & $1750$ \\
    \ce{Bi} & 1.52 & 7.287 & {$+3$，$+5$} & {银白色或微显红色金属} & 9.80 & $271.3$ & $1560$ \\
  \end{tblr}
\end{sidewaystable}

从氮族元素的原子结构来看，它们原子的最外层电子排布都是 $ns^2np^3$，所以，它们在最高氧化物里的化合价是 $+5$ 价。
氮族元素随着原子核外电子层数的增加，获得电子的趋势逐渐减弱，失去电子的趋势逐渐增强，所以，它们的非金属性逐渐减弱，金属性逐渐增强。
氮、磷表现出比较显著的非金属性，砷虽然是非金属，但已表现出一些金属性，而锑、铋已表现出比较明显的金属性。

氮族元素的一些重要性质见\cref{tab:2-2}。

从氮族元素在周期表中的位置来看，氮族元素的非金属性要比同周期的氧族和卤族元素弱。

本章我们主要学习氮和磷。

\begin{Practice}
\begin{question}
  \item 氮族元素包括哪几种元素？分别写出它们的电子排布式。
  \item 氮族元素的性质如何随着电子层数的增加而改变？
  \item 对比磷酸与硝酸、砷酸（\ce{H3AsO4} 一种弱酸）的酸性; 对比磷酸与硫酸、高氯酸（\ce{HClO4}）的酸性。
\end{question}
\end{Practice}

\section{氮气}
氮是一种重要元素，它以双原子分子存在于大气中，约占大气总体积的 78\% 和质量的 75\%。
除了大气是氮的贮藏库外，氮也以化合态形式存在于很多无机物（如硝酸钠和硝酸钾）和有机物（如蛋白质和核酸，二者都是形成生命的重要物质）中。
工业上所用的氮气，通常是以空气为原料，将空气液化后，利用液态空气中液态氮的沸点比液态氧的沸点低而加以分离制得。

\subsection{氮气的物理性质}
纯净的氮气是一种没有颜色、没有气味的气体，比空气稍轻（在标准状况下，\qty{1}{L} 氮气的质量为 \qty{1.25}{g}）。

氮气在 \qty{1}{atm}、 \qty{-195.8}{\celsius} 时，变成没有颜色的液体，\qty{-209.86}{\celsius} 时，变成雪状的固体。

氮气在水里的溶解度很小，在通常状况下，1 体积水中大约可溶解 0.02 体积的氮气。

\subsection{氮气的化学性质}
\par\noindent
\begin{minipage}{0.6\linewidth}\parindent2em
氮分子是由两个氮原子共用三对电子结合而成的分子（\cref{fig:2-1}），氮分子中有三个共价键：
\[\electformularr{N}{d}{d}{d}{d}{d}\:\electformularl{N}{d}{d}{d}{d}{d}
\qquad \ce{N#N}\]
它的键能很大（\qty{226.3}{kCal/mol}），大于其它双原子分子（如氢分子的键能为 \qty{104.2}{kCal/mol}，氧分子的键能为 \qty{118.86}{kCal/mol}），因而氮分子的结构很稳定。
在通常情况下，氮气的性质很不活泼，很难跟其它物质发生化学反应。
但在高温下，氮分子获得了足够的能量，促使它的共价键断裂，就能跟氢、 氧、金属等物质发生化学反应。
\end{minipage}
\begin{minipage}{0.38\linewidth}
\begin{figurehere}
  \includegraphics{2-1.pdf}
  \caption{氮分子示意图}\label{fig:2-1}
\end{figurehere}
\end{minipage}
\par\medskip

\subsubsection{氮气跟氢气的反应}
氮气跟氢气在高温、高压并有催化剂存在的条件下，可以直接化合生成氨（\ce{NH3}）。
\[ \ce{N2 + 3H2 <=>[\text{高温、高压}][\text{催化剂}] 2NH3} + \qty{22.08}{kCal} \]
工业上就是利用这个反应原理来合成氨的。

\subsubsection{氮气跟氧气的反应}
在放电条件下，氮气跟氧气能直接化合生成无色的一氧化氮（\ce{NO}）。
\[ \ce{N2 + O2 \xlongequal{\text{放电}} 2NO} \]

一氧化氮不溶于水，在常温下很容易跟空气中的氧气化合，生成棕色并有刺激性气味的二氧化氮（\ce{NO2}）。
\[ \ce{2NO + O2 \xlongequal{\quad} 2NO2} \]

因此在闪电时，大气中常有少量的二氧化氮产生。

二氧化氮有毒，易溶于水。它溶于水后生成硝酸和一氧化氮。
\[ \ce{3NO2 + H2O \xlongequal{\quad} 2HNO3 + NO} \]

二氧化氮还可相互化合成无色的四氧化二氮（\ce{N2O4}）气体。
\[\schemestart
\chemname{\ce{2NO2}}{\footnotesize （棕色）} 
\arrow(.mid east--.mid west){<=>}
\chemname{\ce{N2O4}}{\footnotesize （无色）} 
\schemestop\chemnameinit{}\]

一氧化氮和二氧化氮是两种重要的氮的氧化物。
氮的氧化物除这两种外，还有一氧化二氮（\ce{N2O}）、三氧化二氮（\ce{N2O3}）、 四氧化二氮（\ce{N2O4}）、五氧化二氮（\ce{N2O5}）等，可见氮跟氧化合时，在不同条件下，它能生成不同的氧化物。
氮在 \ce{N2O}、\ce{NO}、\ce{N2O3}、\ce{NO2}、\ce{N2O5}、这五种不同的氧化物里，其化合价分别是 $+1$、$+2$、$+3$、$+4$、$+5$ 价，氮的最高化合价是 $+5$ 价。

\subsubsection{氮气跟某些金属的反应}
在高温的时候，氮气能够跟镁、钙、锶、钡等金属化合。
例如，镁在空气里燃烧时，除跟氧气化合生成氧化镁外，也能跟氮气化合生成微量的氮化镁（\ce{Mg3N2}）。
\[ \ce{3Mg + N2 \xlongequal{\text{点燃}} Mg3N2}  \]

\subsection{氮的固定}
把大气中的氮转化为氮的化合物叫做氮的固定。
上文所述，在闪电时大气中有氮的氧化物生成，这是自然界的一种氮的固定过程。
豆科作物的根部常附有小根瘤，其中含有能使空气中氮气在常温常压下转化为硝酸盐的固氮菌。
这种转化是自然界的又一种氮的固定过程。

工业上用氮气合成氨，以及在放电条件下制备氮的氧化物，再合成硝酸，这是人工氮的固定过程。

氮化合物的天然矿藏是很少的，而燃料的燃烧、有机物的腐败，又常使氮化合物分解生成氮气进入大气中。
面临对氮的化合物的庞大需要（氮肥、炸药、染料、合成纤维等都需用氮的化合物作原料），当前研究人工氮的固定是十分重要的。

\subsection{氮气的用途}
大量的氮气在工业上用作合成氨的原料，从氨又可以制备各种氮肥和硝酸。
因此，氮气也是一种重要的化工原料。
由于氮气的化学性质不活泼，氮气可用来代替惰性气体作焊接金属时的保护气; 氮气或氮气和氩气的混和气体可用来充填灯泡，以防止钨丝的氧化和减慢钨丝的挥发，使灯泡经久耐用。
粮食、水果如处于低氧高氮的环境中，能使害虫缺氧窒息而死，同时能使植物种子处于休眠状态，代谢缓慢，所以可利用氮气保存粮食、水果等农副产品。

\begin{Practice}
\begin{question}
  \item 为什么在常温下氮气不容易跟其它物质起反应而在高温下却能起反应？
  \item 使含有硫化氢和水蒸气的空气，依次通过氢氧化钠溶液、浓硫酸和灼热的铜丝，最后所得的气体含有哪些成分？ 为什么？写出有关反应的化学方程式。
  \item 五个集气瓶里分别装有下列各种气体: 氯气、氧气、 氮气、二氧化碳和二氧化硫。根据哪些性质鉴别它们: 写出有关反应的化学方程式。
  \item 怎样鉴别二氧化氮和溴蒸气？
\end{question}
\end{Practice}

\section{氨\texorpdfstring{\quad}{ }铵盐}
\subsection{氨}
\subsubsection{氨分子的结构}
\par\medskip\noindent
\begin{minipage}{0.6\linewidth}\parindent2em
氮原子的最外电子层有 5 个电子，电子排布式是 $2s^22p^3$，在氨分子中，氮原子的 3 个未成对的 $2p$ 电子分别与 3 个氢原子中的 $1s$ 电子形成共用电子对，也就是说，氮以三个共价键与三个氢原子联结：
\[ 
\begin{tikzpicture}[baseline=(N.base),inner sep=2.5pt]
  \node (N) {\ce{N}};
  \foreach \x in {1,2,3}
  {
    \node(h\x) at (90*\x:1.2em) {H};
  }
  \pic at ([yshift=-0.4ex]N.west){d}; 
  \pic at ([yshift= 0.4ex]N.west){x}; 
  \pic at ([xshift=-0.4ex]N.north){d}; 
  \pic at ([xshift= 0.4ex]N.north){x}; 
  \pic at ([yshift= 0.4ex]N.east){d}; 
  \pic at ([yshift=-0.4ex]N.east){d}; 
  \pic at ([xshift= 0.4ex]N.south){d}; 
  \pic at ([xshift=-0.4ex]N.south){x};
\end{tikzpicture}
% \begin{tikzpicture}[baseline]
%   \node (c){N};
%   \foreach \x in {1,2,3}
%   {
%     \node(h\x) at (90*\x:1.2em) {H};
%     \draw[draw=none] (c)--(h\x)node[midway,sloped,above,inner sep=1pt]{\scalebox{0.35}{$\times$}}node[midway,sloped,below,inner sep=1pt]{\footnotesize$\cdot$};
%   }
%   \node(h2) at(0:1.2em){\phantom{H}};
%   \draw[draw=none](c)--(h2)node[midway,above,inner sep=0pt]{\footnotesize$\cdot$}node[midway,below,inner sep=1pt]{\footnotesize$\cdot$};
% \end{tikzpicture}
\qquad 
\chemfig{H-N(-[:-90]H)(-[:90]H)} \]

\end{minipage}
\begin{minipage}{0.38\linewidth}
\begin{figurehere}
  \includegraphics{2-2.pdf}
  \caption{氨分子结构示意图}\label{fig:2-2}
\end{figurehere}
\end{minipage}

\medskip
氨分子中的 \ce{N-H} 键是极性键。 
经实验测定，氨分子的结构呈三角锥形，氮原子位于锥顶，三个氢原子位于锥底，\ce{N-H} 键之间的键角为 \ang{107;18;}，所以说氨分子是一个极性分子（\cref{fig:2-2}）。

\subsubsection{氨的物理性质}
氨是没有颜色、具有刺激性气味的气体。
在标准状况下，氨的密度是 \qty{0.771}{g/L}，比同体积的空气轻。

由于氨分子中氢原子是与电负性很大而原子半径较小的氮原子以共价键结合的，因而如同水分子一样，氨分子之间也可形成氢键。
由于氢键的形成，氨分子间的引力增强，使氨很容易液化。
在常压下冷却到 \qty{-33.5}{\celsius} 或在常温下加压到 \qtyrange{7}{8}{atm}，气态氨就凝结为无色的液体，同时放出大量的热。
液态氨气化时要吸收大量的热，而使它周围温度急剧降低，因此，氨常用作致冷剂。

氨极易溶解于水。
在常温下，1 体积水约可溶解 700 体积氨，氨的水溶液叫做氨水。

\subsubsection{氨的化学性质}
\paragraph{氨跟水的反应}\mbox{}\par
\begin{Experiment}*[righthand ratio=0.7]
  在干燥的圆底烧瓶里充满氨气，用带有玻璃管和滴管（滴管里预先吸入水）的塞子塞紧瓶口。立即倒置烧瓶，使玻璃管插入盛有水的烧杯（水里事先加入少量酚酞溶液），挤压滴管的胶头，使少量水进入烧瓶。烧杯里的水即由玻璃管喷入烧瓶，形成红色的喷泉（\cref{fig:2-3}）。
  \tcblower
  \begin{figurehere}
    \includegraphics{2-3.pdf}
    \caption{氨易溶于水}\label{fig:2-3}
  \end{figurehere}
\end{Experiment}
实验表明，氨极易溶于水，大量氨溶于水后，烧瓶里的压强迅速减小，这时外界的大气压就把烧杯里的水迅速压上去。
形成了喷泉。
同时，可以看到这含有少量酚酞的水变成红色。表明氨水具有碱性。

根据加热浓氨水时有氨气逸出，以及氨水在低温时能析出氨的结晶水合物——一水合氨（\ce{NH3.H2O}）的实验事实，可知氨溶于水中，大部分与水结合 成一水合氨（\ce{NH3.H2O}），\ce{NH3.H2O} 是由氨分子和水分子通过氢键结合起来的。

一水合氨很不稳定，受热就会分解而生成氨和水：
\[ \ce{ NH3.H2O \xlongequal{\triangle} NH3 ^ + H2O } \]

一水合氨可以小部分电离成 \ce{NH4+} 和 \ce{OH-}，所以氨水显弱碱性，能使酚酞溶液变红色。
氨在水中的反应可用下式表示：
\[ \ce{ NH3 + H2O <=> NH3.H2O <=> NH4+ + OH-} \]
也可简单表示如下:
\[ \ce{ NH3 + H2O <=> NH4+ + OH-} \]

\paragraph{氨跟酸的反应}\mbox{}\par
\begin{Experiment}*[righthand ratio=0.5]
  拿一根玻璃棒在浓氨水里蘸一下，另一根玻璃棒在浓盐酸里蘸一下，使这两根玻璃棒接近（不要接触），就有大量的白烟产生（\cref{fig:2-4}）。
  \tcblower
  \begin{figurehere}
    \includegraphics{2-4.pdf}
    \caption{氨跟氯化氢反应}\label{fig:2-4}
  \end{figurehere}
\end{Experiment}

这白烟是氨水里挥发出的氨跟浓盐酸挥发出的氯化氢化合所生成的微小的氯化铵晶体：
\[ \ce{NH3 + HCl \xlongequal{\quad} NH4Cl} \]

氨同样能跟其它的酸化合生成铵盐。
如把氨通入硝酸或硫酸中，就会生成硝酸铵或硫酸铵：
\begin{gather*}
  \ce{ NH3 + HNO3 \xlongequal{\quad} NH4NO3 }\\
  \ce{ 2NH3 + H2SO4 \xlongequal{\quad} (NH4)2SO4 }
\end{gather*}

\paragraph{氨跟氧气的反应}\mbox{}\par
在催化剂（如铂、氧化铁、氧化铬等）存在的情况下，氨跟氧气发生如下的反应：
\[ \ce{4NH3 + 5O2 }\xlongequal[\triangle]{\text{催化剂}} \ce{4NO + 6H2O} + \qty{216.7}{kCal} \]

\begin{Experiment}*[righthand ratio=0.28]
慢慢把空气通入盛浓氨水的锥形瓶中，再将红热的螺旋状铂丝\footnote{用加热 \ce{(NH4)2Cr2O7} 新制得的 \ce{Cr2O3} 代替铂丝作催化剂，实验效果也很明显。}接近液面，但不要使铂丝跟氨水接触（\cref{fig:2-5}）。可观察到锥形瓶中有棕色气体生成\footnote{由于锥形瓶里有水蒸气存在，它跟二氧化氮起反应生成硝酸，再跟氨起反应生成硝酸铵晶体。因此，锥形瓶里常看到有白烟产生。}，这是因为氨被氧化为一氧化氮，后者再遇氧气生成二氧化氮。同时，还可观察到铂丝继续保持红热，这是因为氨分子跟氧分子在铂丝表面上进行的这个反应是放热的。
\tcblower 
\begin{figurehere}
  \includegraphics{2-5.pdf}
\caption{氨的催化氧化}\label{fig:2-5}
\end{figurehere}
\end{Experiment}

上述反应叫做氨的催化氧化（或叫接触氧化），它是工业上制硝酸的基础。

\subsubsection{氨的实验室制法}
在实验室里常用给铵盐和碱加热的方法来制取氨。
\[ \ce{ 2NH4Cl + Ca(OH)2 \xlongequal{\triangle} CaCl2 + 2NH3 ^ + 2H2O} \]
\begin{Experiment}*[righthand ratio=0.7]
给试管里的氯化铵和消石灰的混和物加热，用倒立的干燥的试管收集氨（\cref{fig:2-6}）。把润湿的红色石蕊试纸放在试管口，观察试纸颜色的变化，可以试验氨是否已经充满试管。
\tcblower 
\begin{figurehere}
  \includegraphics{2-6.pdf}
\caption{氨的制取}\label{fig:2-6}
\end{figurehere}
\end{Experiment}

实验室中要制取干燥的氨，通常使制得的氨通过碱石灰\footnote{在氢氧化钠浓溶液中加入氧化钙，加热，制成白色固体即得碱石灰，它是水分和二氧化碳的吸收剂。}，以吸收其中的水蒸气。

\subsubsection{氨的用途}
氨是一种重要的化工产品。
它不仅是氮肥工业的基础，同时又是制造硝酸、铵盐、纯碱等的重要原料。 
氨在有机合成工业（如合成纤维、塑料、染料等）里也是一种常用的原料。 氨又是一种常用的致冷剂。



\subsection{铵盐}
氨跟酸作用可生成铵盐。
铵盐是由铵离子（\ce{NH4+}）和酸根离子组成的化合物。

铵盐都是晶体，能溶解于水。铵盐的化学性质如下：
\subsubsection{铵盐受热分解}
\begin{Experiment}*[righthand ratio=0.6]
  给试管里的氯化铵晶体加热，观察发生的现象（\cref{fig:2-7}）。
  \tcblower
  \begin{figurehere}
    \includegraphics{2-7.pdf}
    \caption{氯化铵受热分解}\label{fig:2-7}
  \end{figurehere}
\end{Experiment}

受热时，氯化铵分解生成氨和氯化氢，冷却时，它们又重新结合生成氯化铵。
\begin{gather*}
  \ce{ NH4Cl \xlongequal{\triangle} NH3 ^ + HCl ^ }\\
  \ce{ NH3 + HCl \xlongequal{\quad} NH4Cl   }
\end{gather*}

碳酸氢铵受热时，分解生成氨、水和二氧化碳。
\[ \ce{ NH4HCO3 \xlongequal{\triangle} NH3 ^ + H2O +CO2 ^} \]
铵盐受热容易分解，分解时，一般放出氨。

\subsubsection{铵盐跟碱的反应}
铵盐能跟碱起反应放出氨气。例如：
\[ \ce{ (NH4)2SO4 +2NaOH \xlongequal{\triangle} Na2SO4 +2NH3 ^ + 2H2O } \]

这个性质是一切铵盐的通性。
实验室里就利用这样的反应来制取氨，同时也可以利用这个性质来检验铵离子的存在。

铵盐在工农业生产上有着重要的用途。
大量的铵盐用作氮肥。
硝酸铵还用来制炸药。
氯化铵常用作印染和制干电池的原料，它也用在金属的焊接上，以除去金属表面上的氧化物薄层。




\begin{Practice}
\begin{question}
  \item 浓硫酸常用作气体的干燥剂，是否可用来干燥氨气？为什么？
  \item 碘受热变成蒸气，碘蒸气遇冷变成碘；氯化铵受热分解生成的气体遇冷仍变成氯化铵，这两种现象本质上是否相同？为什么？
  \item 怎样用化学方法证明硫酸铵既是铵盐又是硫酸盐？写出检验方法、现象和有关的化学方程式。
  \item 制备 \qty{1}{L} 含氨 10\% 的氨水（密度是 \qty{0.96}{g/cm^3}）需要用多少体积的氨气（标准状况下）？
  \item 350 体积（标准状况下）的氨溶解在 1 体积的水里，这种氨水的百分比浓度是多少？摩尔浓度是多少？（这种氨水的密度是 \qty{0.924}{g/cm^3}）
  \item 用氢氧化钙和氯化铵各 \qty{10}{g}，在标准状况下可以制得多少升的氨气？如果把这些氨配成 \qty{500}{ml} 的氨水，这种溶液的摩尔浓度是多少？
\end{question}
\end{Practice}


\section{硝酸的工业制法}
硝酸是一种重要的化工产品，它是制造炸药、染料、塑料、硝酸盐和许多其它化工产品的重要原料。

现代生产硝酸最重要的方法是氨的催化氧化法，这个方法的生产过程大致可分为两个阶段：
\begin{enumerate*}[(1)]
  \item 氨氧化生成一氧化氮；
  \item 一氧化氮氧化生成二氧化氮，二氧化氮被水（或稀硝酸）吸收而生成硝酸。
\end{enumerate*}
\subsubsection{氨的氧化}
\par\medskip\noindent
\begin{minipage}{0.6\linewidth}\parindent2em
将氨和净化后的空气以一定比混和，进入氧化炉（\cref{fig:2-8}）。
氧化炉的中部有多层水平的铂铑合金网作为催化剂，在 \qty{800}{\celsius} 的高温下，氨跟氧气在网上进行反应生成一氧化氮和水蒸气，同时放出大量的热。
\[ \ce{ 4NH3 + 5O2} \xlongequal[\text{高温}]{\ce{Pt}-\ce{Rh}} \ce{4NO + 6H2O } + \qty{216.7}{kCal}\]
\end{minipage}\hfill
\begin{minipage}{0.38\linewidth}
\begin{figurehere}
  \includegraphics{2-8.pdf}
  \caption{氧化炉}\label{fig:2-8}
\end{figurehere}
\end{minipage}\par\medskip

\subsubsection{硝酸的生成}
一氧化氮经过冷却，再被空气中的氧氧化成二氧化氮。
\[ \ce{ 2NO + O2 \xlongequal{\quad} 2NO2 } + \qty{27.02}{kCal} \]

最后，在吸收塔内二氧化氮被水吸收，就得到了硝酸。
\[ \ce{ 3NO2 + H2O \xlongequal{\quad} 2HNO3 +NO } + \qty{32.5}{kCal} \]
从这一反应看，只有 2/3 的二氧化氮转化为硝酸，而 1/3 的二氧化氮转化为一氧化氮。
因此，常在吸收反应进行过程中补充一些空气，使生成的一氧化氮再氧化为二氧化氮，二氧化氮溶于水又生成硝酸和一氧化氮。
经过这样多次的氧化和吸收，二氧化氮可以比较完全地被水吸收，能够尽可能多地转化为硝酸。

从吸收塔出来的尾气中尚含有少量未被吸收的一氧化氮和二氧化氮，如果不加处理就排放到空气中，会造成污染，严重危害人体健康及农作物生长。
为消灭氮的氧化物对大气的污染，变废为宝，将尾气通入碱液吸收塔，用碱液吸收，可制得重要的化工原料亚硝酸钠：
\[ \ce{ NO + NO2 +2NaOH \xlongequal{\quad} 2NaNO2 + H2O } \]

用上述方法制得的硝酸的浓度一般为 50\% 左右，如果要制取更浓的硝酸，可用硝酸镁（或浓硫酸）作为吸水剂，将稀硝酸蒸馏浓缩，就可以得到 96\% 以上的浓硝酸。

\begin{Practice}[习题]
\begin{question}
  \item 简述氨氧化制取硝酸的化学反应原理，并写出有关反应的化学方程式。
  \item 雷雨时，雨水里可含有微量的硝酸。人们曾模拟自然界的这一过程，用电弧法生产硝酸。当空气通过电极之间的电弧时，极高的温度，导致氮气氧化为一氧化氮。但即使在 \qty{3000}{\celsius} 的高温，也只能生成 5\%（按体积计）的 \ce{NO}。由于耗电多，产率低，这个方法已逐渐被其它方法代替。试写出用电弧法生产硝酸的有关的化学方程式。
  \item 氨氧化制硝酸时，如果由氨制成一氧化氮的产率是 96\%，由一氧化氮制成硝酸的产率是 92\%。\qty{10}{t} 氨可以制备多少吨 50\% 的硝酸？
  \item 制备 \qty{6}{M} \ce{HNO3} 溶液 \qty{250}{ml}，问需用 65.3\% 的浓硝酸（密度是 \qty{1.4}{g/cm^3}）多少毫升？
\end{question}
\end{Practice}


\section{硝酸\texorpdfstring{\quad}{ }硝酸盐}
\subsection{硝酸}
\subsubsection{硝酸的物理性质}
纯硝酸是无色、易挥发、有刺激性气味的液体，密度为 \qty{1.5027}{g/cm^3}，沸点 \qty{83}{\celsius}，凝固点 \qty{-42}{\celsius}。
它能以任意比溶解于水。
常用的浓硝酸的浓度大约是 69\%。
浓度为 98\% 以上的浓硝酸在空气里由于硝酸的挥发而产生“发烟”现象，通常叫做发烟硝酸。
这是因为硝酸里放出的硝酸蒸气遇到空气里的水蒸气生成了极微小的硝酸液滴的缘故。

\subsubsection{硝酸的化学性质}
硝酸是一种强酸。
它除了具有酸的通性以外，还有它本身的特性。

\paragraph{硝酸的不稳定性}\mbox{}\par
硝酸很不稳定，容易分解。
纯净的硝酸或浓硝酸在常温下见光就会分解，受热时分解得更快。
\[ \ce{ 4HNO3 \xlongequal{\triangle} 2H2O + 4NO2 ^ + O2 ^ }\]

硝酸越浓，就越容易分解。
分解放出的二氧化氮溶于硝酸而使硝酸呈黄色。
为了防止硝酸分解，必须把它盛在棕色瓶里，贮放在黑暗而且温度低的地方。

\paragraph{硝酸的氧化性}\mbox{}\par
硝酸是一种很强的氧化剂，不论稀硝酸还是浓硝酸都有氧化性，几乎能跟所有的金属（除金、铂等少数金属外）或非金属发生氧化—还原反应。
\begin{Experiment}
  在放有铜片的两个试管里，分别加入少量浓硝酸和稀硝酸，观察现象。
\end{Experiment}
浓硝酸和稀硝酸都能跟铜起反应。
前者反应激烈，有红棕色的气体产生；后者反应较缓慢，有无色气体产生，在试管口变红棕色。

以上反应的化学方程式分别是：
\begin{gather*}
  \ce{ Cu + 4HNO3\text{（浓）}\xlongequal{\quad} Cu(NO3)2 + 2NO2 ^ + 2H2O }\\
  \ce{ 3Cu + 8HNO3\text{（稀）}\xlongequal{\quad} 3Cu(NO3)2 + 2NO ^ +4H2O }
\end{gather*}

从上述两个反应可以看出: 硝酸跟金属发生反应时，主要是正五价的氮得到电子，被还原成较低价的氮的化合物，并不象盐酸跟较活泼金属反应那样放出氢气。
除金、铂等少数几种金属外，硝酸几乎可以使所有的金属氧化而生成硝酸盐。

值得注意的是，有些金属如铝、铁等常温下在浓硝酸中会发生钝化现象。
这是因为浓硝酸将它们的表面氧化成一层薄而致密的氧化物薄膜，因而阻止了进一步反应的缘故。
所以，可以用铝槽车装盛浓硝酸。

浓硝酸和浓盐酸的混和物（摩尔比 $1:3$）叫做王水，它的氧化能力更强，能使一些不溶于硝酸的金属，如金、铂等溶解。

硝酸还能使许多非金属（如碳、硫、磷）及某些有机物（如松节油、锯末等）氧化。例如：
\[ \ce{ 4HNO3 + C \xlongequal{\quad} 2H2O +4NO2 ^ +CO2 ^ }\]


\subsubsection{硝酸的实验室制法}
硝酸有挥发性，所以，在实验室里可以把硝酸盐跟浓硫酸共同加热来制取它。
\[ \ce{ NaNO3 + H2SO4\text{（浓）}\xlongequal{\quad} NaHSO4 + HNO3  ^ } \]

\subsection{硝酸盐}
多数硝酸盐是无色晶体，极易溶于水。
硝酸盐性质不稳定，加热易分解放出氧气，所以，在高温时硝酸盐是强氧化剂。
\begin{Experiment}
给试管里的硝酸钾加热到熔融，把带有火星的细木条伸进试管口检验放出的气体，同时观察放出气体的颜色和试管中剩余物的颜色。

分别用硝酸铜、硝酸银晶体代替硝酸钾，做上面的实验。
\end{Experiment}
从实验可知，各种硝酸盐加热后都会分解，一般地说，在金属活动性顺序表里镁以前的比较活泼金属的硝酸盐，加热时放出氧气并生成亚硝酸盐。例如：
\[ \ce{ 2KNO3 \xlongequal{\triangle} 2KNO2 + O2 ^ }\]

在金属活动性顺序表里位于镁和铜之间的金属的硝酸盐，在加热时，生成金属氧化物、二氧化氮和氧气。例如：
\[ \ce{ 2Cu(NO3)2 \xlongequal{\triangle} 2CuO +4NO2 ^ + O2 ^ }\]

活动性很小的金属的硝酸盐，即在金属活动性顺序表里铜以后的金属的硝酸盐，在加热时生成金属单质、二氧化氮和氧气。例如：
\[ \ce{ 2AgNO3 \xlongequal{\triangle} 2Ag + 2NO2 ^ + O2 ^ } \]

由于硝酸盐是氧化剂，所以，用硝酸钾和易燃的硫粉、炭粉可配制成黑色火药。

\begin{Practice}[习题]
\begin{question}
  \item 硝酸和硫酸、盐酸的性质有什么相同的地方和不同的地方？怎样用实验方法来鉴别这三种酸？
  \item 依次写出下列变化的化学方程式，并注明反应发生的条件。
  \[\ce{N2}\to\ce{NH3}\to\ce{NO}\to\ce{NO2}\to\ce{HNO3}\to\ce{NH4NO3}\]
  \item 把铜片放到下列各种酸里各有什么现象发生？ 能起反应的写出化学方程式，不能起反应的说明理由。
  \begin{tasks}(5)
    \task 稀盐酸，
    \task 浓硫酸，
    \task 稀硫酸，
    \task 浓硝酸，
    \task 稀硝酸。
  \end{tasks}
  \item \qty{3.5}{M} 的稀硝酸溶液 \qty{50}{ml} 与足量的铜起反应，生成多少克硝酸铜？计算生成的一氧化氮在标准状况下的体积。
  \item 写出下列四种硝酸盐受热分解的化学方程式：
  \begin{tasks}(4)
    \task \ce{NaNO3}；
    \task \ce{Ca(NO3)2}；
    \task \ce{Fe(NO3)2}；
    \task \ce{Hg(NO3)2}。
  \end{tasks}
  \item 62\% 的硝酸溶液（密度为 \qty{1.38}{g/cm^3}）的摩尔浓度是多少？若将这种浓硝酸 \qty{100}{ml} 稀释至 \qty{500}{ml}，求它的摩尔浓度是多少？
  \item 给硝酸钠晶体跟浓硫酸和铜共热，将会发生什么现象？用这种方法可以鉴定硝酸盐吗？对于很稀的硝酸盐溶液用这种方法鉴定行不行？
\end{question}
\end{Practice}

\section{氧化—还原反应方程式的配平}
在初中化学里，我们学过一些氧化—还原反应，这些氧化—还原反应的化学方程式比较简单，反应物和生成物的系数是较小的整数，不难通过观察来配平。
但是，在高中化学里已经学到的或以后将会学到的较复杂的氧化—还原反应，如硝酸跟金属或非金属的反应等，就不容易用观察的方法来配平它们的化学方程式。
那么，这些复杂的化学方程式怎样配平呢？

我们知道，氧化—还原反应的本质是参加反应的原子间的电子转移（包括电子得失和电子对的偏移），原子间的电子转移可以用元素的化合价的升降来表示。
因此，氧化—还原反应的化学方程式可以通过分析电子转移或化合价升降来配平。
在这里，我们学习用化合价升降的方法来配平化学方程式。

\begin{example}
  配平铜跟稀硝酸反应的化学方程式。
\end{example}
\begin{solution}
  \begin{enumerate}[1.]
    \item 先写出反应物和生成物的化学式，并列出发生氧化和还原的元素的正负化合价。
    \[ \OX{a,\ox*{0,Cu}} + \ce{H}\OX{b,\ox*{+5,N}}\ce{O_3} \ce{ -> } \OX{c,\ox*{+2,Cu}}\ce{(NO3)2} + \OX{d,\ox*{+2,N}}\ce{O} + \ce{H2O}
\]
    \item 列出元素的化合价的变化。
    \par\vspace{5mm}
    \[
  \OX{a,\ox*{0,Cu}} + \ce{H}\OX{b,\ox*{+5,N}}\ce{O_3} \ce{ -> } \OX{c,\ox*{+2,Cu}}\ce{(NO3)2} + \OX{d,\ox*{+2,N}}\ce{O} + \ce{H2O}
  \redox(a,c)[-stealth]{\scriptsize 化合价升高 2}
  \redox(b,d)[-stealth][-1]{\scriptsize 化合价降低 3}
\]\par\medskip\mbox{}
    \item 使化合价的升高和降低的总数相等。
    \par\vspace{5mm}
    \[
  3\,\OX{a,\ox*{0,Cu}} + \ce{2H}\OX{b,\ox*{+5,N}}\ce{O_3} \ce{ -> } 3\,\OX{c,\ox*{+2,Cu}}\ce{(NO3)2} + 2\,\OX{d,\ox*{+2,N}}\ce{O} + \ce{H2O}
  \redox(a,c)[-stealth]{\scriptsize 化合价升高 $2\times3$}
  \redox(b,d)[-stealth][-1]{\scriptsize 化合价降低 $3\times2$}
\]\par\medskip\mbox{}
    \item 用观察的方法配平其它物质的系数。在上述反应里，有 6 个 \ce{NO3-} 没有参与氧化—还原反应，所以 \ce{HNO3} 的系数应是 8；\ce{H2O} 的系数应是 4，因为有 2 个 \ce{NO3-} 还原成 \ce{NO}，其中 4 个氧原子跟 \ce{HNO3} 中氢离子结合成水。配平后，把单线改成等号。
    \[ \ce{ 3Cu + 8HNO3\,\text{（稀）} \xlongequal{\quad} 3Cu(NO3)2 + 2NO ^ +4H2O } \]
  \end{enumerate}
\end{solution}

\begin{example}
  配平碳跟硝酸起反应的化学方程式。
\end{example}
\begin{solution}
  \begin{enumerate}[1.]
    \item 写出反应物和生成物的化学式，列出发生氧化—还原反应的元素的正负化合价。
    \[
  \OX{a,\ox*{0,C}} + \ce{H}\OX{b,\ox*{+5,N}}\ce{O_3} \ce{ -> } \OX{c,\ox*{+4,N}}\ce{O2} + \OX{d,\ox*{+4,C}}\ce{O2} + \ce{H2O}
  \]
    \item 列出元素的化合价变化。
    \par\vspace{5mm}
    \[
  \OX{a,\ox*{0,C}} + \ce{H}\OX{b,\ox*{+5,N}}\ce{O_3} \ce{ -> } \OX{c,\ox*{+4,N}}\ce{O2} + \OX{d,\ox*{+4,C}}\ce{O2} + \ce{H2O}
  \redox(a,d)[-stealth]{\scriptsize 化合价升高 4}
  \redox(b,c)[-stealth][-1]{\scriptsize 化合价降低 1}
\]\par\medskip\mbox{}
    \item 使化合价的升高和降低的总数相等。
    \par\vspace{5mm}
    \[
  \OX{a,\ox*{0,C}} + 4\,\ce{H}\OX{b,\ox*{+5,N}}\ce{O_3} \ce{ -> } 4\,\OX{c,\ox*{+4,N}}\ce{O2} + \OX{d,\ox*{+4,C}}\ce{O2} + \ce{H2O}
  \redox(a,d)[-stealth]{\scriptsize 化合价升高 4}
  \redox(b,c)[-stealth][-1]{\scriptsize 化合价降低 $1\times4$}
\]\par\medskip\mbox{}
    \item 配平其它物质的系数，把单线改成等号。
    \[ \ce{ C + 4HNO3 \xlongequal{\quad} 4NO2 ^ + CO2 ^ + 2H2O } \]
  \end{enumerate}
\end{solution}

\begin{Practice}[习题]
\begin{question}
  \item 配平下列氧化—还原反应的化学方程式，并指出哪种元素被氧化了，哪种元素被还原了。
  \begin{tasks}
    \task \ce{ SO2 + O2 -> SO3 }
    \task \ce{ P + Cl2 -> PCl3 }
    \task \ce{ NaBr + Cl2 -> NaCl + Br2 }
    \task \ce{ Fe + FeCl3 -> FeCl2 }
    \task \ce{ KClO3 -> KCl + O2 }
  \end{tasks}
  \item 配平下列氧化—还原反应的化学方程式。
  \begin{tasks}
    \task \ce{ Cu + H2SO4 \text{（浓）} -> CuSO4 + SO2 + H2O }
    \task \ce{ MnO2 + HCl \text{（浓）} -> MnCl2 + Cl2 + H2O }
    \task \ce{ NH3 + O2 -> NO + H2O }
    \task \ce{ NO2 + H2O -> HNO3 + NO }
    \task \ce{ Mg + HNO3 \text{（稀）} -> Mg(NO3)2 + NH4NO3 + H2O}
  \end{tasks}
  \item 配平过氧化钠跟二氧化碳的反应。
  \[ \ce{ Na2O2 + CO2 -> Na2CO3 + O2 } \]

  （提示: 过氧化钠中氧元素的化合价是 $-1$，在反应中，一个氧原子从 $-1$ 降为 $-2$，另一个氧原子从 $-1$ 升为 0。）
\end{question}
\end{Practice}

\section{磷\texorpdfstring{\quad}{ }磷酸\texorpdfstring{\quad}{ }磷酸盐}
\subsection{磷}
由同种元素组成的不同性质的单质叫做同素异形体。
磷有多种同素异形体，其中最重要的是白磷和红磷。

\subsubsection{磷的物理性质}
白磷是一种蜡状的固体，有剧毒，不溶于水，但能溶于二硫化碳里。
把白磷隔绝空气加热到 \qty{260}{\celsius}，就会转变成红磷。
红磷是暗红色粉末状的固体，没有毒，不溶于水，也不溶于二硫化碳。
红磷加热到 \qty{416}{\celsius} 时就升华，它的蒸气冷却后变成白磷。

\subsubsection{磷的化学性质}
磷的化学性质活泼，容易跟氧、卤素以及许多金属直接化合。
\paragraph{磷跟氧的化合反应}\mbox{}\par
\begin{Experiment}*[righthand ratio=0.65]
  把一块铁片水平地夹在铁架上，把少量的白磷和红磷隔开相当的距离分放在铁片上（\cref{fig:2-9}），然后在红磷的下面加热，观察是红磷还是白磷先着火燃烧。
  \tcblower
  \begin{figurehere}
    \includegraphics{2-9.pdf}
    \caption{白磷和红磷的着火点的比较}\label{fig:2-9}
  \end{figurehere}
\end{Experiment}

实验说明白磷远比红磷容易燃烧。白磷的着火点是 \qty{40}{\celsius}，红磷的着火点是 \qty{240}{\celsius}。
白磷受到轻微的摩擦或被加热到 \qty{40}{\celsius}，就会发生燃烧现象。
所以，白磷必须贮存在密闭容器里，少量的白磷可保存在水里。

白磷和红磷的着火点虽然不同，但是燃烧以后，都生成五氧化二磷。
五氧化二磷极易吸水，是一种强干燥剂。

白磷在空气里，即使在常温下，也会缓慢地氧化，氧化时会发光，在暗处可以清楚地看见。

磷跟氮一样，在它的氧化物和氯化物里的最高化合价是 $+5$，在它跟氢和金属的化合物里，化合价是 $-3$。

\paragraph{磷跟卤素化合}\mbox{}\par
由于卤素的电负性比磷高，因此，磷在其卤化物中显示 $+3$ 和 $+5$ 价。
磷在不充足的氯气中燃烧生成三氯化磷。
\[ \ce{2P + 3Cl2 \xlongequal{\text{点燃}} 2PCl3} \]

在过量的氯气中燃烧生成五氯化磷。
\[ \ce{2P + 5Cl2 \xlongequal{\text{点燃}} 2PCl5} \]

\par\medskip\noindent
\begin{minipage}{0.55\linewidth}\parindent2em
白磷和红磷在性质上的差别，是由于它们具有不同的结构引起的。
白磷分子是由四个磷原子结合而成的 \ce{P4} 分子，4 个磷原子排布在一个正四面体的 4 个顶角上，每个磷原子与分子中的其他 3 个磷原子以共价键相结合（\cref{fig:2-10}）。而红磷的结构远比白磷复杂，这里就不作介绍了。
\end{minipage}\hfill
\begin{minipage}{0.4\linewidth}
\begin{figurehere}
  \includegraphics{2-10.pdf}
  \caption{白磷分子结构示意图}\label{fig:2-10}
\end{figurehere}
\end{minipage}

\subsubsection{磷的存在和用途}
磷在空气中易被氧化，因此自然界里没有游离态的磷存在。
磷主要以磷酸盐的形式存在于矿石中。
此外动物的骨骼、牙齿、脑髓和神经组织里都含有磷。
植物的果实和幼芽里也含有磷。
磷对于维持生物体正常的生活机能有重要的作用。

白磷可用于制造纯度高的磷酸; 红磷除用于制农药外，主要用于制造安全火柴。
火柴盒侧面所涂的物质就是红磷和三硫化二锑等的混和物，而火柴头上的物质一般是氧化剂（氯酸钾、二氧化锰）和易燃物如硫等。
当两者摩擦时，因摩擦而生的热使跟氯酸钾接触的红磷发火，并引起火柴头上的易燃物燃烧，从而使火柴杆着火。
此外，在军事上还用磷来制造烟幕弹和燃烧弹。

\subsection{磷酸和磷酸盐}
五氧化二磷极易与水化合，发生剧烈反应，同时放出大量的热。
随着反应条件的不同，可生成偏磷酸\footnote{从一分子磷酸脱去一分子水而成的酸，称为偏磷酸。\ce{H3PO4 = HPO3 + H2O}}（\ce{HPO3}）或磷酸（\ce{H3PO4}）：
\begin{gather*}
  \ce{ P2O5 +  H2O \xlongequal{\text{冷水}}  2HPO3 }\\
  \ce{ P2O5 + 3H2O \xlongequal{\text{热水}} 2H3PO4 }
\end{gather*}

磷酸是无色透明的晶体，熔点为 \qty{42.35}{\celsius}，具有吸湿性，易溶于水，和水能以任何比混和。通常用的磷酸是一种无色粘稠的浓溶液，内含 83\%～98\% 的纯磷酸。

磷酸没有毒，而偏磷酸则有剧毒。

磷酸比硝酸稳定，不易分解。
工业上是用硫酸跟磷酸钙反应来制取磷酸的：
\[ \ce{ Ca3(PO4)2 +3H2SO4 \xlongequal{\quad} 2H3PO4 + 3CaSO4 v } \]

滤去硫酸钙沉淀，所得滤液就是磷酸溶液。

磷酸不显氧化性，是一种中等强度的三元酸。
它能形成三种类型的盐，一种正盐和两种酸式盐。
例如：

\begin{tblr}{colspec={*4{l}},hlines=0pt,vlines=0pt}
  磷酸二氢盐 & \ce{NaH2PO4} & \ce{Ca(H2PO4)2}  & \ce{NH4H2PO4}   \\
  磷酸氢盐   & \ce{Na2HPO4} & \ce{CaHPO4}      & \ce{(NH4)2HPO4} \\
  磷酸盐     & \ce{Na3PO4}  & \ce{Ca3(PO4)2}   & \ce{(NH4)3PO4}  \\
\end{tblr}

所有的磷酸二氢盐都易溶于水，而磷酸氢盐和磷酸盐中除钾、钠和铵盐外，几乎都不溶于水。

磷酸盐大量用于磷肥。
自然界的磷矿石的主要成分是磷酸钙，它是难溶于水的矿物。
化学工业上制造磷肥的目的，就是加工磷矿石，使它们由难溶于水的正盐转化为较易溶于水（或弱酸）的酸式盐，以利于植物吸收。

\begin{Practice}[习题]
\begin{question}
  \item 白磷和红磷在性质上有何不同？怎样证明白磷和红磷是同素异形体？白磷是哪种类型的晶体？
  \item 怎样分别用下列两种物质为原料来制备磷酸：
  \begin{enumerate*}[(1)]
    \item 磷，
    \item 磷酸钙。
  \end{enumerate*}
  如果要制备 \qty{500}{g} 磷酸，需要磷、磷酸钙各多少克？制备过程中的化学反应是否都属于氧化—还原反应？
  \item 现有 $2.4M$ \ce{H3PO4} 溶液 \qty{500}{ml}，需要分别加入多少毫升 $3M$ \ce{NaOH} 溶液才能把反应物中所含 \ce{H3PO4} 全部转化为：
  \begin{enumerate*} [(1)]
    \item 磷酸二氢钠，
    \item 磷酸氢二钠，
    \item 磷酸钠。
  \end{enumerate*}
  \item 当 \qty{26.4}{g} 的硫酸铵跟过量的氢氧化钠一起加热时，放出的气体全部被含 \qty{39.2}{g} 磷酸的溶液所吸收。这时生成的盐是什么？
  \item 写出下列物质依次变化的化学方程式：
  \[ \ce{P}\to\ce{P2O5}\to\ce{H3PO4}\to\ce{Ca3(PO4)2}\to\ce{Ca(H2PO4)2}\]
\end{question}
\end{Practice}

\section*{内容提要}
\addcontentsline{toc}{section}{内容提要}\setcounter{subsection}{0}
\subsection*{一、氮族元素}
氮族元素属于元素周期表的第 \MyRoman{5} 主族，包括氮、磷、 砷、锑、铋五种元素。
氮族元素原子的最外电子层的排布为 $ns^2np^3$，其非金属性比同周期的氧族、卤族元素弱。
\subsection*{二、氮气}
\begin{enumerate}[1.]
  \item 由于氮分子（\electformularr{N}{d}{d}{d}{d}{d}\:\electformularl{N}{d}{d}{d}{d}{d}）的键能很大，结构稳定，在常温下它很不活泼。但在高温下，氮分子也能跟氢气、氧气、金属等许多物质起化学反应。
  \item 氮有很多氧化物，其中一氧化氮和二氧化氮是两种重要的氮的氧化物。
\end{enumerate}
\subsection*{三、氮和铵盐}
\begin{enumerate}[1.]
  \item 氨分子的结构呈三角锥形，它是极性分子。由于氨分子间可形成氢键，所以氨容易液化; 由于氨分子与水分子间可形成氢键，所以氨极易溶于水。
  \item 氨水具有弱碱性，氨溶于水发生下列反应:
  \[ \ce{NH3 + H2O <=> NH3.H2O <=> NH4+ +OH-} \]
  \item 氨跟酸反应生成铵盐。
  \item 铵盐受热易分解。
  \item 铵盐跟碱起反应：
  \[ \ce{ NH4+ OH- \xlongequal{\triangle} NH3 ^ + H2O } \]
  实验室利用这个反应来检验铵离子（\ce{NH4+}）的存在。
\end{enumerate}
\subsection*{四、硝酸和硝酸盐}
\begin{enumerate}[1.]
  \item 硝酸除具有酸的通性外，还具有以下特性：
  \begin{enumerate}[(1)]
    \item 不稳定性。见光受热易分解。
    \item 氧化性。硝酸是一种很强的氧化剂，它几乎能跟所有的金属（除金、铂外）、非金属发生氧化—还原反应。
  \end{enumerate}
  \item 氨氧化法制硝酸的化学反应大致可分为两个阶段，即氮催化氧化生成一氧化氮; 一氧化氮氧化生成二氧化氮，二氧化氮被水吸收生成硝酸。
  \item 硝酸盐性质不稳定，受热易分解放出氧气，所以在高温下也是强氧化剂。
\end{enumerate}

\subsection*{五、磷、磷酸和磷酸盐}
\begin{enumerate}[1.]
  \item 磷有多种同素异形体，其中最重要的是白磷和红磷。
  \item 磷能够跟氧气、卤素等直接化合。
  \item 磷酸是一种中等强度的三元酸，比硝酸稳定，不易分解。
  \item 磷酸能形成三种类型的盐：正盐、氢盐和二氢盐。
\end{enumerate}
\subsection*{六、用化合价升降的方法来配平氧化—还原反应方程式}
先写出反应物和生成物的化学式，并列出发生氧化或还原的元素的正负化合价，再列出元素的化合价变化并使化合价的升高和降低的总数相等，最后用观察的方法配平其它物质的系数。

\begin{Review}
\begin{question}
  \item 解释下列现象：
  \begin{tasks}
    \task 为什么浓硝酸常放在棕色的瓶子里？
    \task 锌跟稀硫酸反应可产生氢气，而跟稀硝酸反应则没有氢气放出。
  \end{tasks}
  \item 有一种白色晶体，它跟氢氧化钠共热的时候，放出一种无色气体，这种气体能使润湿的红色石蕊试纸变蓝; 跟浓硫酸共热的时候，也放出一种无色气体，这种气体能使润湿的蓝色石蕊试纸变红。如果这两种气体相遇会发生白烟。原来的白色晶体可能是什么物质？写出上述有关反应的化学方程式。
  \item 怎样利用空气和水为原料来制造硝酸铵？写出制造过程中有关反应的化学方程式并注明反应条件。
  \item 硝酸铜可以用下列三种方法制备：铜跟浓硝酸起反应，铜跟稀硝酸起反应，氧化铜跟硝酸起反应。如果用这三种方法制取等量的硝酸铜，消耗的纯硝酸的质量是否相等？哪一种用掉的硝酸最少？
  \item 根据下列氮元素化合价的变化，写出实现各步变化的化学方程式，并分别指出 (1)、(4) 两步变化中什么物质是还原剂？
  \[\schemestart 
    \chemfig{ \charge{90[anchor=-90]={\tiny 0}}{N} }
    \arrow{->[][\small (1)]}
    \chemfig{ \charge{90[anchor=-90]={\tiny $-3$}}{N} }
    \arrow{->[][\small (2)]}
    \chemfig{ \charge{90[anchor=-90]={\tiny $+2$}}{N} }
    \arrow{->[][\small (3)]}
    \chemfig{ \charge{90[anchor=-90]={\tiny $+4$}}{N} }
    \arrow{->[][\small (4)]}
    \chemfig{ \charge{90[anchor=-90]={\tiny $+5$}}{N} }
  \schemestop\]
  \item 氯化钠和氯化铵在性质上有什么相似的地方？怎样从氯化铵和氯化钠的混和物里把氯化铵分离出来？
  \item 有五瓶白色固体: 硝酸钠、硫酸铵、碳酸钙、氯化铵、磷酸钠，你用什么化学方法把它们分别检验出来？写出实验操作步骤和产生的现象，以及它们的化学方程式。
  \item 某一种气态氮的氧化物 \qty{250}{ml}（在标准状况下），质量为 \qty{0.33}{g}，它的组成里含氧 53.5\%，求它的化学式。
  \item 取铜银合金 \qty{1}{g}，放入硝酸中溶解，再加盐酸，则得氯化银沉淀 \qty{0.35}{g}，求此合金中铜、银的百分率各多少？
  \item 把 \qty{10}{ml} 一氧化氮和二氧化氮的混和气体通入倒立在水槽的盛满水的量筒里，片刻以后，量筒里留下 \qty{5}{ml} 气体。计算通入的混和气体里，一氧化氮和二氧化氮各几毫升。
  \item 配平下列氧化—还原反应的化学方程式：
  \begin{tasks}
    \task \ce{ KMnO4 + HCl -> MnCl2 + KCl + H2O + Cl2 }
    \task \ce{ FeS2 + O2 -> Fe2O3 +SO2 }
  \end{tasks}
\end{question}
\end{Review}